\begin{abstract}
University support offices can hold only a limited number of outreach meetings and must decide
early which students to contact. This study asks whether such an office should rely on a black-box
model for this decision, taking the tabular foundation model \tabicl as the test case, and whether
the model's explanations are trustworthy and useful to people. Using 3,630 Dropout or Graduate
records from the public UCI \emph{Predict Students' Dropout and Academic Success} dataset, all
modelling choices were fixed on a validation split before a single locked final test, in which each
of five models flagged the top 20\% of students, the assumed outreach capacity. \tabicl was not
reliably better than XGBoost (ROC-AUC difference 0.001, 95\% interval $-0.006$ to 0.009) or an
interpretable Explainable Boosting Machine (EBM; 0.003, $-0.007$ to 0.014): all three reached a
ROC-AUC between 0.954 and 0.957, and at least 108 of each model's 109 flagged students dropped out.
A post-hoc audit of all 544 validation students showed that the SHAP explanations were faithful to
the model but often unstable: when the comparison group changed, 411 students fell below the 0.80
stability threshold for their top five reasons. The flags also largely restated information already
visible to advisers, since 78 of the 108 flagged validation students had passed no first-semester
units. In a 50-person forward-simulation study, there was no evidence that explanation charts
improved participants' predictions of the model's flag (0.900 with charts versus 0.916 without;
difference $-0.016$, 95\% interval $-0.060$ to 0.032). On this dataset, the added complexity of the
foundation model is not justified. The interpretable EBM, whose explanations are exact rather than
estimated, is the stronger candidate for further validation, and any explanation shown to staff
should state its uncertainty and support, rather than replace, an adviser's judgement.
\end{abstract}

\noindent\textbf{Keywords:} student dropout risk, explainable artificial intelligence,
tabular foundation models, interpretable machine learning, SHAP, human-centred evaluation
